\section{Public and private are two corners, not two definitions}
\label{sec:framework}

Take the network as the unit of analysis: agents are nodes, a request from one to another
is an edge. An edge is not one thing. It is a triple
%
\[
  e = (F,\; S,\; A)
\]
%
--- \arch{$F$ormation}, who may open it; \arch{$S$emantics}, what it carries once open;
\txt{$A$ttributes}, what is asserted about it. The three are independently settable, and
the first two are checked by a coordination layer on every call while the third only ever
reaches a model as text.

\begin{quote}
\textbf{Public and private are not points on one dial, and they are not definitions we
supply. They are two characteristic \emph{bundles} in this space} --- two corners that
co-occur in deployed systems because each component makes the others affordable. Every
real network sits somewhere between them, and most sit at no corner at all.
\end{quote}

Twitter and WhatsApp are the two corners. On Twitter an account addresses any other
without assent, relevance is a claim in a profile, exchanges reveal nothing but their
content, and trust arrives as a platform-computed badge; a bad actor changes handles. On
WhatsApp a conversation exists because a number was already held or someone made an
introduction, the tie carries files and history, what is known about the counterpart is
who introduced them and how the last exchange went, and misbehaviour costs the tie and
embarrasses whoever vouched.

The consequence is the paper's premise:

\begin{quote}
\textbf{``Public versus private'' is not an experimental variable.} It is a comparison
between two bundles that differ on three axes at once. A study that runs it measures
$(F_\text{pub}, S_\text{pub}, A_\text{pub}) \rightarrow (F_\text{pri}, S_\text{pri},
A_\text{pri})$ and reports one number, which cannot be attributed to any component. The
variables are $do(F)$, $do(S)$ and $do(A)$, one at a time.
\end{quote}

\ifcompact\else
\begin{figure}[t]
\centering
\begin{tikzpicture}[font=\small,>=Latex,
  ag/.style={circle,draw=black!45,fill=black!3,minimum size=5.5mm,inner sep=0pt},
  req/.style={circle,draw=black,fill=black!10,minimum size=6.5mm,inner sep=0pt,font=\scriptsize\bfseries},
  ttl/.style={font=\footnotesize\bfseries,align=center}]

% ---- panel 1: formation
\begin{scope}
  \node[ttl,arch] at (1.55,2.05) {edge formation};
  \node[font=\scriptsize,gray!70,align=center] at (1.55,1.62) {which edges may exist};
  \node[req] (R1) at (0,0.55) {R};
  \node[ag] (a1) at (1.3,1.15) {}; \node[ag] (a2) at (1.3,-0.05) {};
  \draw[->,black!65] (R1) -- (a1); \draw[->,black!65] (R1) -- (a2);
  \draw[dotted,black!45,rounded corners=6pt] (0.85,-0.55) rectangle (1.75,1.6);
  \node[ag] (b1) at (2.95,1.35) {}; \node[ag] (b2) at (2.95,0.55) {}; \node[ag] (b3) at (2.95,-0.25) {};
  \draw[->,dashed,attr] (R1) to[out=25,in=170] (b1);
  \draw[->,dashed,attr] (R1) to[out=-25,in=190] (b3);
  \node[font=\scriptsize,attr] at (2.05,-0.9) {assent or broker};
\end{scope}

% ---- panel 2: semantics
\begin{scope}[xshift=5.05cm]
  \node[ttl,arch] at (1.35,2.05) {edge semantics};
  \node[font=\scriptsize,gray!70,align=center] at (1.35,1.62) {what an existing edge carries};
  \node[req] (R2) at (0,0.55) {R};
  \node[ag] (c1) at (2.7,1.25) {}; \node[ag] (c2) at (2.7,-0.15) {};
  \draw[->,black!70] (R2) to[out=30,in=180] node[midway,above,font=\scriptsize,black]{task} (c1);
  \draw[<-,black!70] ([yshift=-2.2mm]R2.east) to[out=8,in=200] node[midway,below,font=\scriptsize,black,yshift=1pt]{result} ([xshift=-1.5mm,yshift=-2.6mm]c1.west);
  \draw[<->,black!70] (R2) to[out=-30,in=180] node[midway,below,font=\scriptsize,black]{list + read} (c2);
  \draw[black!25,fill=black!6,rounded corners=1pt] (3.05,-0.42) rectangle (3.55,0.12);
  \node[font=\scriptsize,gray!70] at (3.3,-0.75) {store};
\end{scope}

% ---- panel 3: attributes
\begin{scope}[xshift=9.9cm]
  \node[ttl,attr] at (1.35,2.05) {edge attributes};
  \node[font=\scriptsize,gray!70,align=center] at (1.35,1.62) {what is asserted about it};
  \node[req] (R3) at (0,0.55) {R};
  \node[ag] (d1) at (2.7,0.55) {};
  \draw[->,black!70] (R3) -- (d1);
  \node[draw=attr,fill=attr!7,rounded corners=1.5pt,inner sep=2.5pt,font=\scriptsize,align=center]
    at (1.35,1.15) {\textit{``trusted, 2 years''}};
  \node[font=\scriptsize,gray!70,align=center] at (1.35,-0.25) {reaches the model\\[-2pt]as tokens only};
\end{scope}
\end{tikzpicture}
\caption{The three axes, drawn on $G=(V,E)$. \arch{Formation} and \arch{semantics} are
enforced: a coordination layer refuses an edge that may not form and a request an edge does
not carry. \txt{Attributes} are never checked by anything---they arrive in a prompt, which
is why asking whether a model acts on them is a real question rather than a tautology. Each
regime supplies a characteristic value on all three (Table~\ref{tab:axes}), and a study that
compares public with private moves all three at once.}
\label{fig:axes}
\end{figure}
\fi

Table~\ref{tab:axes} is what each corner supplies on each axis --- the bundles unpacked.

\begin{table}[t]
\centering\small
\caption{The three axes, and what each regime supplies on it. Formation and semantics
are enforced by a coordination layer on every call; attributes only ever reach a model
as text. A study that compares a public network to a private one moves all three.}
\label{tab:axes}
\begin{tabular}{@{}p{0.155\linewidth}p{0.375\linewidth}p{0.375\linewidth}@{}}
\toprule
& \textbf{Public} & \textbf{Private} \\
\midrule
\arch{Edge\newline formation}
& Unilateral. Anyone in the directory is addressable; matching is by declared
  relevance. \textbf{Lateral topology}: wide and shallow, many weak ties.
& Bilateral or brokered. A counterpart is reachable because of a prior tie or an
  introduction. \textbf{Longitudinal topology}: narrow and deep, few strong ties. \\
\addlinespace[2pt]
\arch{Edge\newline semantics}
& Sandboxed by default. A counterpart accepts a task and returns a result; its store
  stays opaque. Interaction is one-shot request/response.
& Openable. Shared workspaces, persistent namespaces, a store that may be listed and
  read---affordable precisely because the tie exists. \\
\addlinespace[2pt]
\txt{Edge\newline attributes}
& Platform-computed for strangers: reputation scores, ratings, usage counts,
  verification badges.
& Carried by the relationship: tenure, who introduced whom, how the last
  collaboration went, who is answerable. \\
\bottomrule
\end{tabular}
\end{table}

\ifcompact\else Figure~\ref{fig:axes} draws the three on the graph.\fi The rest of this
section takes them one at a time.

\ifcompact
\paragraph{The two topologies are the substance of the formation axis.} A public directory
buys \emph{breadth}: a rare requirement finds someone and---less obviously---a claim can be
put to a second source. A private roster buys \emph{depth}: fewer counterparts, each
already paid for, each with a history. These are not two settings of one dial; they are the
returns on two different edge prices. Which links form, and at what price, is strategic
network formation \citep{jackson1996}, and the agent setting differs in that the price is
conventionally zero. On \arch{semantics}, exposure follows the tie---one does not open a
working directory to a stranger---but the axis is separately settable, and a capability
kernel expresses both settings as grants over different namespaces, so a condition is a set
of capabilities rather than a different codebase \citep{a2a2025}. On \txt{attributes}, both
regimes carry a characteristic form of trust-conveying metadata that no runtime checks: a
public network computes a score about a stranger, a private one carries a history about an
acquaintance. Both rest on the same assumption---that an agent builds a representation of
its counterpart from what it is told and acts differently because of it.

\else
\paragraph{Edge formation} is which edges a requester may create: the scope of the
directory it can enumerate and address. This is the axis the word ``public'' usually
refers to, and the one broadcast networks for agents are built around---an agent
publishes what it needs and what it offers, and the network routes to whoever is
relevant \citep{eigenflux}. Which links form, and at what price, is strategic network
formation \citep{jackson1996}; the agent setting differs in that the price is
conventionally zero.

The two topologies are the substance of the axis. A public directory buys
\emph{breadth}: a rare requirement finds someone, and---less obviously---a claim can be
put to a second source. A private roster buys \emph{depth}: fewer counterparts, each
already paid for, each with a history. Breadth and depth are not two settings of one
dial; they are the returns on two different edge prices.

\paragraph{Edge semantics} is what an existing edge carries. A counterpart may accept a
task and return a result while revealing nothing else, or expose a store to be listed
and read. The distinction is not ours: it separates agent-to-agent task protocols, which
deliberately keep a counterpart opaque \citep{a2a2025}, from granted resource scopes, in
which one principal hands another read access to files. It maps onto the regimes because
exposure follows the tie---one does not open a working directory to a stranger---but it
is separately settable, and a capability kernel expresses both as grants over different
namespaces, so a condition is a set of capabilities rather than a different codebase.

\paragraph{Edge attributes} is what is asserted \emph{about} an edge. Both regimes have
a characteristic form of it, and both are trust-conveying metadata that no runtime
checks: a public network computes a score about a stranger, a private one carries a
history about an acquaintance. They rest on the same assumption---that an agent builds
a representation of its counterpart from what it is told, and acts differently because
of it.

\fi

\subsection{Why these three and not nine}

The literature on ties, teams and open innovation supplies at least nine relational
dimensions---membership openness, identity continuity, shared-state depth, repeated
interaction, relationship origin, prior trust, shared goals, accountability, exit cost.
They are not nine degrees of freedom: in a deployed system an open directory brings weak
identity, which brings diffuse accountability, which brings low exit cost. The three-way
grouping breaks that coupling the one way a controlled study can. Formation and semantics
are properties of the interaction layer---independently settable, checked on every call,
not open to argument by a model. Attributes are properties of the people behind the agents;
a harness cannot instantiate them, only \emph{assert} them. That asymmetry is the
experiment rather than a limitation to apologise for: deployed networks communicate trust,
alignment, attribution and tenure as text on cards, so asking whether a model acts on that
text asks whether a large part of agent-network design does anything at all.

\ifcompact
\subsection{Two corrections the framework has to carry}

Both fall out of taking the edge, rather than visibility, as the primitive. First,
\textbf{a public network's noise is partly a feature}: a requester who can address many
counterparts can put the same question to a second one, and contradiction is
\emph{detectable}, so a roster that filters misinformation out before it arrives also
removes the second source needed to catch what does arrive. Whether the filter or the
redundancy dominates is empirical, and the framework must not settle it by calling breadth
a cost. Second, \textbf{an endorsement is not an introduction}, and they sit on different
axes. \txt{``This counterpart is reliable''} is an attribute: text on an edge that already
exists. \arch{An introduction that makes a third agent addressable} is formation: it
changes the graph. A framework that files both under ``trust'' predicts that a private
network's advantage comes from what it knows about its members; the decomposition predicts
something narrower---that only the half which changes reachability can do work---and
\S\ref{sec:evidence} reports which way it went.
\else
\subsection{Two corrections the framework has to carry}

Both fall out of taking the edge, rather than visibility, as the primitive.

\paragraph{A public network's noise is partly a feature.} The usual complaint is that an
open directory is full of unvetted claims. It is---but a requester who can address many
counterparts can put the same question to a second one, and contradiction is
\emph{detectable}. A roster that filters misinformation out before it arrives also
removes the second source needed to catch what does arrive. Whether the filter or the
redundancy dominates is an empirical question, not a definitional one, and the framework
must not settle it by calling breadth a cost.

\paragraph{An endorsement is not an introduction.} These are on different axes and are
routinely conflated. \txt{``This counterpart is reliable''} is an attribute: text, on an
edge that already exists. \arch{An introduction that makes a third agent addressable} is
formation: it changes the graph. A framework that files both under ``trust'' predicts
that a private network's advantage comes from what it knows about its members. The
decomposition predicts something narrower---that only the half which changes reachability
can do work---and \S\ref{sec:evidence} reports which way it went.
\fi
